\textbf{Phần 1. Trắc nghiệm lựa chọn}
\setcounterpageref{ex}{0}

\begin{ex} %[3D16N78]
Kết quả của phép nhân $30\,000 \times 3$ là
\choice
{\True $90000$}
{$60000$}
{$100000$}
{$80000$}
\loigiai{
$30000 \times 3 = 90000$.\\
\textbf{Đáp án của câu hỏi là A}
}
\end{ex}

\begin{ex} %[3D16N78]
Kết quả của phép chia $80\,000 : 4$ là
\choice
{$15000$}
{$25000$}
{\True $20000$}
{$30000$}
\loigiai{
$80000 : 4 = 20000$.\\
\textbf{Đáp án của câu hỏi là C}
}
\end{ex}

\begin{ex} %[3D16H78]
Kết quả của phép nhân $15\,210 \times 4$ là
\choice
{$60830$}
{\True $60840$}
{$60850$}
{$59840$}
\loigiai{
$15210 \times 4 = 60840$.\\
\textbf{Đáp án của câu hỏi là B}
}
\end{ex}

\begin{ex} %[3D16H78]
Kết quả của phép chia $69\,369 : 3$ là
\choice
{$23122$}
{$23124$}
{\True $23123$}
{$22123$}
\loigiai{
$69369 : 3 = 23123$.\\
\textbf{Đáp án của câu hỏi là C}
}
\end{ex}

\begin{ex} %[3D16H78]
Tính giá trị biểu thức: $12\,000 \times 5 + 25\,000$.
\choice
{\True $85000$}
{$80000$}
{$90000$}
{$75000$}
\loigiai{
$60000 + 25000 = 85000$.\\
\textbf{Đáp án của câu hỏi là A}
}
\end{ex}

\begin{ex} %[3D16H78]
Tính giá trị biểu thức: $(80\,000 - 20\,000) : 3$.
\choice
{$15000$}
{\True $20000$}
{$25000$}
{$30000$}
\loigiai{
$60000 : 3 = 20000$.\\
\textbf{Đáp án của câu hỏi là B}
}
\end{ex}

\begin{ex} %[3D16V78]
Một nhà máy sản xuất 4 ngày, mỗi ngày được $16\,500$ sản phẩm. Đã xuất đi $45\,000$ sản phẩm. Số sản phẩm còn lại là
\choice
{$20000$ sản phẩm}
{$22000$ sản phẩm}
{$25000$ sản phẩm}
{\True $21000$ sản phẩm}
\loigiai{
Tổng: $16500 \times 4 = 66000$ sản phẩm.\\
Còn lại: $66000 - 45000 = 21000$ sản phẩm.\\
\textbf{Đáp án của câu hỏi là D}
}
\end{ex}

\begin{ex} %[3D16V78]
Một khu đất hình vuông có cạnh dài $14\,000\text{ m}$. Chu vi khu đất đó là
\choice
{\True $56000\text{ m}$}
{$52000\text{ m}$}
{$60000\text{ m}$}
{$48000\text{ m}$}
\loigiai{
$14000 \times 4 = 56000\text{ m}$.\\
\textbf{Đáp án của câu hỏi là A}
}
\end{ex}

\begin{ex} %[3D16V78]
Tìm $x$ biết: $x : 6 = 15\,000$.
\choice
{$80000$}
{\True $90000$}
{$100000$}
{$85000$}
\loigiai{
$x = 15000 \times 6 = 90000$.\\
\textbf{Đáp án của câu hỏi là B}
}
\end{ex}

\begin{ex} %[3D16V78]
Tìm $y$ biết: $y \times 4 = 84\,880$.
\choice
{$21210$}
{$21230$}
{$21240$}
{\True $21220$}
\loigiai{
$y = 84880 : 4 = 21220$.\\
\textbf{Đáp án của câu hỏi là D}
}
\end{ex}

\begin{ex} %[3D16V78]
Một cuộn dây cáp dài $75\,000\text{ m}$ cắt thành 5 đoạn bằng nhau. Mỗi đoạn dài là
\choice
{$14000\text{ m}$}
{$16000\text{ m}$}
{\True $15000\text{ m}$}
{$18000\text{ m}$}
\loigiai{
$75000 : 5 = 15000\text{ m}$.\\
\textbf{Đáp án của câu hỏi là C}
}
\end{ex}

\begin{ex} %[3D16C78]
Có hai kho hàng chứa tất cả $90\,000\text{ kg}$ gạo. Kho thứ nhất chứa $35\,000\text{ kg}$. Hỏi kho thứ hai chứa nhiều hơn kho thứ nhất bao nhiêu ki-lô-gam gạo?
\choice
{$18000\text{ kg}$}
{$22000\text{ kg}$}
{$25000\text{ kg}$}
{\True $20000\text{ kg}$}
\loigiai{
Kho 2: $90000 - 35000 = 55000\text{ kg}$.\\
Kho 2 hơn kho 1: $55000 - 35000 = 20000\text{ kg}$.\\
\textbf{Đáp án của câu hỏi là D}
}
\end{ex}

\textbf{Phần 2. Trắc nghiệm đúng sai}
\setcounterpageref{ex}{0}

\begin{ex} %[3D16H78]
Xét tính đúng sai của các phép tính:
\choiceTF
{\True $25000 \times 4 = 100000$.}
{\True $96000 : 3 = 32000$.}
{$15000 \times 6 = 80000$.}
{\True $15000 \times 6 = 90000$.}
\loigiai{
\begin{itemize}
    \item $25000 \times 4 = 100000$: Đúng.
    \item $96000 : 3 = 32000$: Đúng.
    \item $15000 \times 6 = 90000 \neq 80000$. Ý 3 Sai, Ý 4 Đúng.
\end{itemize}
\begin{itemize}
    \item a) \textbf{Đúng}.
    \item b) \textbf{Đúng}.
    \item c) \textbf{Sai}.
    \item d) \textbf{Đúng}.
\end{itemize}
}
\end{ex}

\begin{ex} %[3D16V78]
Một kho gạo có $84\,000\text{ kg}$. Lần 1 xuất $\dfrac{1}{3}$ số gạo, lần 2 xuất $26\,000\text{ kg}$. Xét tính đúng sai:
\choiceTF
{\True Lần 1 xuất được $28\,000\text{ kg}$ gạo.}
{\True Cả 2 lần xuất được $54\,000\text{ kg}$ gạo.}
{Kho còn lại $32\,000\text{ kg}$ gạo.}
{\True Kho còn lại $30\,000\text{ kg}$ gạo ($84000 - 54000 = 30000$).}
\loigiai{
\begin{itemize}
    \item Lần 1: $84000 : 3 = 28000\text{ kg}$. Đúng.
    \item Tổng xuất: $28000 + 26000 = 54000\text{ kg}$. Đúng.
    \item Còn lại: $84000 - 54000 = 30000\text{ kg} \neq 32000$. Ý 3 Sai, Ý 4 Đúng.
\end{itemize}
\begin{itemize}
    \item a) \textbf{Đúng}.
    \item b) \textbf{Đúng}.
    \item c) \textbf{Sai}.
    \item d) \textbf{Đúng}.
\end{itemize}
}
\end{ex}

\textbf{Phần 3. Trả lời ngắn}
\setcounterpageref{ex}{0}

\begin{ex} %[3D16H78]
Tính tích: $18\,000 \times 5$.
\shortans{90000}
\loigiai{
\textbf{Đáp số:} $90000$\\
$18000 \times 5 = 90000$.
}
\end{ex}

\begin{ex} %[3D16H78]
Tính thương: $75\,000 : 3$.
\shortans{25000}
\loigiai{
\textbf{Đáp số:} $25000$\\
$75000 : 3 = 25000$.
}
\end{ex}

\begin{ex} %[3D16V78]
Tính giá trị biểu thức: $14\,000 \times 5 - 30\,000$.
\shortans{40000}
\loigiai{
\textbf{Đáp số:} $40000$\\
$70000 - 30000 = 40000$.
}
\end{ex}

\begin{ex} %[3D16V78]
Tính giá trị biểu thức: $96\,000 : 4 + 14\,000$.
\shortans{38000}
\loigiai{
\textbf{Đáp số:} $38000$\\
$24000 + 14000 = 38000$.
}
\end{ex}

\begin{ex} %[3D16V78]
Tìm $x$ biết: $x \times 5 = 85\,000$.
\shortans{17000}
\loigiai{
\textbf{Đáp số:} $17000$\\
$x = 85000 : 5 = 17000$.
}
\end{ex}

\begin{ex} %[3D16C78]
Một hình chữ nhật có chu vi $80\,000\text{ m}$, chiều dài gấp 3 lần chiều rộng. Tính chiều dài hình chữ nhật theo mét.
\shortans{30000}
\loigiai{
\textbf{Đáp số:} $30000$\\
Nửa chu vi: $80000 : 2 = 40000\text{ m}$.\\
Chiều dài: $40000 : 4 \times 3 = 30000\text{ m}$.
}
\end{ex}
